\appendix
\setcounter{equation}{0}
\renewcommand{\theequation}{A.\arabic{equation}}
\renewcommand{\theHequation}{A.\arabic{equation}}

\section{Proofs and Demand Details}
\label{app:proofs}

\begin{proof}[Proof of Lemma~\ref{lem:demand}]
Write \(\gamma=1-\alpha-\beta>0\). The first-order conditions for the
outer household problem are
\[
 \frac{\alpha}{M}=\lambda P_M,\qquad
 \frac{\beta}{N}=\lambda,\qquad
 \frac{\gamma}{Z}=\lambda.
\]
The budget constraint gives \(\lambda=1/Y\), hence expenditures on
\((M,N,Z)\) equal \((\alpha Y,\beta Y,\gamma Y)\). Substitution into utility
gives
\begin{equation}
 R=K\frac{Y}{P_M^\alpha},\qquad
 K=\alpha^\alpha\beta^\beta\gamma^\gamma.
 \label{eq:indirect-utility-proof}
\end{equation}

Let \(h_P(p_A,p_P,1)\) be the platform Hicksian demand for one unit of the
homothetic retail aggregate. Differentiability of the unit-expenditure
function and Shephard's lemma imply
\[
 \frac{\partial P_M}{\partial p_P}=h_P(p_A,p_P,1),\qquad
 s=\frac{p_Ph_P(p_A,p_P,1)}{P_M}
   =\frac{\partial\ln P_M}{\partial\ln p_P}.
\]
Because \(p_A\) is fixed and \(d\ln p_P/dz=-1\), the chain rule yields
\(d\ln P_M/dz=-s\), proving the demand identities in the lemma.
\end{proof}

\begin{proof}[Proof of Lemma~\ref{lem:equilibrium}]
For fixed \((z,\kappa,G)\), define
\[
 T(Y)=B+\omega G+(\beta+\alpha\omega)Y.
\]
Since \(s\in(0,1)\) and both retention shares lie in \([0,1]\),
\(0\leq\omega\leq1\). Thus \(T\) maps \(\mathbb R_+\) into itself and
\[
 |T(Y)-T(\widetilde Y)|
 = (\beta+\alpha\omega)|Y-\widetilde Y|,
 \qquad \beta+\alpha\omega\leq\alpha+\beta<1.
\]
The contraction theorem gives a unique fixed point and convergence of
iteration from every nonnegative initial value. Solving the fixed-point
equation yields \eqref{eq:income-solution}. Its numerator is positive, while
\(D\geq1-\alpha-\beta>0\), so \(Y>0\). Lemma~\ref{lem:demand} then determines
the remaining household objects uniquely.
\end{proof}

\subsection{Benchmark multipliers}

Holding \((z,\kappa,G)\) fixed and varying the realized autonomous-income
injection or outside-funded retail spending gives
\[
 \frac{\partial Y}{\partial B}=\frac1D,
 \qquad
 \frac{\partial Y}{\partial G}=\frac{\omega}{D}
 =\omega\sum_{n=0}^{\infty}(\beta+\alpha\omega)^n.
\]
The series converges because \(\beta+\alpha\omega<1\). It records that retail
leakage recurs in every induced spending round.

\begin{proof}[Proof of Proposition~\ref{prop:finite-capacity}]
The ratio identities follow immediately from \(Y=B/D\) and
\(R=KY/P_M^\alpha\). Organizational capacity does not affect \(B_j\),
\(s_j\), or \(P_{Mj}\). Differentiating
\[
 \mathcal D(\kappa)=\frac{D_1(\kappa)}{D_0(\kappa)}
\]
therefore gives
\begin{align}
 \frac{\partial\ln\mathcal D(\kappa)}{\partial\kappa}
 &=-\alpha\delta_\rho
 \left[\frac{s_1}{D_1(\kappa)}-
       \frac{s_0}{D_0(\kappa)}\right] \notag\\
 &=-\frac{\alpha\delta_\rho d_A(s_1-s_0)}
 {D_1(\kappa)D_0(\kappa)}<0.
 \label{eq:finite-capacity-monotonicity-proof}
\end{align}
Hence both finite income ratios are strictly increasing in \(\kappa\).

For any \(q\in(\mathcal D(1),\mathcal D(0))\), continuity and strict
monotonicity give a unique \(\kappa(q)\in(0,1)\) satisfying
\(\mathcal D(\kappa(q))=q\). Solving
\[
 d_A+\alpha s_1\Delta(\kappa)
 =q[d_A+\alpha s_0\Delta(\kappa)]
\]
gives \eqref{eq:finite-capacity-solution}. The denominator in that expression
is positive at a solution because
\[
 s_1-qs_0
 =\frac{s_1D_0(\kappa)-s_0D_1(\kappa)}{D_0(\kappa)}
 =\frac{d_A(s_1-s_0)}{D_0(\kappa)}>0.
\]
Since \(0<\mathcal P<1\),
\(q_R^F=\mathcal B/\mathcal P^\alpha>q_Y^F=\mathcal B\). The inverse of the
strictly decreasing function \(\mathcal D\) is decreasing, so
\(\kappa_R^F<\kappa_Y^F\). Comparing \(\mathcal D(\kappa)\) with the two
targets proves the three strict regions. At the first threshold,
\[
 \frac{R_1}{R_0}=1,\qquad
 \frac{Y_1}{Y_0}=\mathcal P^\alpha<1;
\]
at the second,
\[
 \frac{Y_1}{Y_0}=1,\qquad
 \frac{R_1}{R_0}=\mathcal P^{-\alpha}>1.
\]
\end{proof}

\begin{proof}[Proof of Corollary~\ref{cor:increasing-differences}]
At \(G=0\),
\[
 Y=\frac{B}{D},\qquad
 D=d_A+\alpha s\Delta(\kappa),\qquad
 D_\kappa=-\alpha s\delta_\rho.
\]
Therefore \(\partial\ln Y/\partial\kappa=\alpha s\delta_\rho/D\). Since
\(D_z=\alpha\Delta(\kappa)s_z\),
\begin{align}
 \frac{\partial^2\ln Y}{\partial z\,\partial\kappa}
 &=\alpha\delta_\rho
 \frac{s_zD-sD_z}{D^2} \\
 &=\frac{\alpha\delta_\rho d_A s_z}{D^2}>0.
\end{align}
Finally, \(\ln R=\ln K+\ln Y-\alpha\ln P_M\), and \(P_M\) is independent
of \(\kappa\), so the same cross-partial holds for \(\ln R\).
\end{proof}

\begin{proof}[Proof of Proposition~\ref{prop:general-share-speed}]
Strict concavity, the endpoint restrictions, and interior positivity imply
\(g'(0)>0\) and \(g'(1)<0\). For example, for every \(u\in(0,1)\),
\[
 g'(0)\geq\frac{g(u)-g(0)}u>0,
 \qquad
 g'(1)\leq\frac{g(1)-g(u)}{1-u}<0.
\]
Write \(b=\alpha\Delta(\kappa)>0\) and define
\[
 F(s,b)=g'(s)(d_A+bs)-bg(s).
\]
Direct differentiation gives
\begin{equation}
 \frac{\partial\Lambda_g}{\partial s}
 =\frac{bF(s,b)}{(d_A+bs)^2},
 \qquad
 F_s(s,b)=g''(s)(d_A+bs)<0.
 \label{eq:general-peak-foc}
\end{equation}
Moreover,
\[
 F(0,b)=d_Ag'(0)>0,
 \qquad
 F(1,b)=(d_A+b)g'(1)<0.
\]
Thus \(F\) has exactly one zero \(s^*(b)\in(0,1)\). Equation
\eqref{eq:general-peak-foc} shows that \(\Lambda_g\) rises before this zero and
falls after it. Since its endpoint values are zero and it is positive in the
interior, \(s^*\) is the unique global maximizer.

If \(g(s)=g(1-s)\), differentiability gives \(g'(1/2)=0\), so
\(F(1/2,b)=-bg(1/2)<0\). Because \(F(0,b)>0\) and \(F\) is strictly
decreasing, \(s^*(b)<1/2\).

For the capacity comparative statics, evaluate expressions at \(s=s^*(b)\)
and write \(D^*=d_A+bs^*\). The first-order condition implies
\(g'(s^*)=bg(s^*)/D^*\), and hence
\[
 F_b(s^*,b)=s^*g'(s^*)-g(s^*)
 =-\frac{d_Ag(s^*)}{D^*}<0.
\]
The implicit-function theorem, \(b_\kappa=-\alpha\delta_\rho\), and the
envelope theorem yield
\begin{align}
 \frac{\partial s^*}{\partial\kappa}
 &=-\frac{\alpha\delta_\rho d_Ag(s^*)}
 {g''(s^*)[D^*]^2}>0, \\
 \frac{\partial\Lambda_g^{\max}}{\partial\kappa}
 &=-\frac{\alpha\delta_\rho d_Ag(s^*)}{[D^*]^2}<0.
 \label{eq:general-peak-comparative-statics-proof}
\end{align}
\end{proof}

\subsection{Admissibility of the general share-speed}

The share-speed restriction is consistent with a two-channel homogeneous demand
system. Let \(t=\ln(p_P/p_A)\) and take any interior solution of
\(ds/dt=-g(s)\). Define
\[
 \phi(r)=\exp\left\{\int_{r_0}^{r}s(\ln u)\,d\ln u\right\},
 \qquad
 P_M(p_A,p_P)=p_A\phi(p_P/p_A).
\]
Then \(P_M\) is linearly homogeneous and its platform expenditure share is
\(p_P P_{M,p_P}/P_M=s\). It is increasing because the two expenditure shares
are \(s\) and \(1-s\). Moreover,
\[
 \phi''(r)
 =-\frac{\phi(r)}{r^2}\left[g(s)+s(1-s)\right]<0,
\]
so \(P_M\) is a valid concave unit-expenditure function. This construction also
shows that the non-CES share-speed in Figure~\ref{fig:mechanisms} is not merely
an arbitrary curve detached from a demand system.

\subsection{CES specialization}
\label{app:ces}

The main assumptions nest the two-channel CES aggregator
\begin{equation}
 M=\left[(1-a)^{1/\eta}M_A^{(\eta-1)/\eta}
 +a^{1/\eta}M_P^{(\eta-1)/\eta}\right]^{\eta/(\eta-1)},
 \qquad a\in(0,1),\quad\eta>1.
 \label{eq:ces-specialization}
\end{equation}
Cost minimization gives
\begin{align}
 P_M(z)&=\left[(1-a)p_A^{1-\eta}
 +ap_P(z)^{1-\eta}\right]^{1/(1-\eta)}, \\
 s(z)&=\frac{ap_P(z)^{1-\eta}}
 {(1-a)p_A^{1-\eta}+ap_P(z)^{1-\eta}}, \\
 s'(z)&=(\eta-1)s(1-s).
 \label{eq:ces-price-share}
\end{align}
Thus \(g(s)=(\eta-1)s(1-s)\) satisfies all conditions in
Proposition~\ref{prop:general-share-speed}, including symmetry. With
\(t(\kappa)=\sqrt{1+\alpha\Delta(\kappa)/d_A}\), its unique peak and height
have the closed forms
\begin{align}
 s^*(\kappa)&=\frac{1}{1+t(\kappa)}<\frac12, \\
 \Lambda_g^{\max}(\kappa)
 &=(\eta-1)\frac{t(\kappa)-1}{t(\kappa)+1}.
 \label{eq:ces-peak-closed-form}
\end{align}
These expressions are the CES specialization, not assumptions needed for the
general peak result.

As \(\eta\downarrow1\), the CES system approaches the Cobb--Douglas channel
case, \(P_M\to p_A^{1-a}p_P^a\) and \(s\to a\). At fixed unequal prices,
\(\eta\to\infty\) concentrates expenditure on the cheaper channel. These are
limits of the specialization rather than finite-parameter corners in the main
interior demand path.

Figure~\ref{fig:local-marginal-regions} combines this CES specialization with
the marginal-income identities in the main text. It is a local derivative
map, not a finite-change calculation: for each fixed capacity level, the
position of \(\Lambda_g(s,\kappa)\) relative to \(\varepsilon\) and
\(\varepsilon+\alpha s\) determines the signs of \(d\ln Y/dz\) and
\(d\ln R/dz\). Movement to the right can represent deeper integration only
conditional on the maintained mapping \(ds/dz>0\) and fixed remaining
parameters. At finite CES prices the share lies strictly in \((0,1)\); the
displayed endpoints are the continuous extension used by the peak theorem.
The CES weight \(a\) and the two price levels determine the mapping from \(z\)
to \(s\), but do not enter this geometry once \(s\) is placed on the horizontal
axis.

\begin{figure}[H]
 \centering
 \includegraphics[width=0.97\textwidth]{supplementary_local_marginal_regions.pdf}
 \caption{Local marginal income regions in the CES specialization. Under
 \(G=0\), the curves plot \(\Lambda_g(s,\kappa)\) at \(\kappa=0,0.5,1\) for
 \(g(s)=3s(1-s)\), \(\alpha=0.35\),
 \(\beta=0.25\), \(\eta=4\), \(\varepsilon=0.10\), \(\rho_A=0.8\), and
 \((\underline\rho_P,\overline\rho_P)=(0.1,0.6)\). The horizontal line is
 \(\varepsilon\); the upward-sloping line is \(\varepsilon+\alpha s\).
 Below \(\varepsilon\), nominal local income and consumption-equivalent real
 income both rise. Between the two reference lines, nominal local income falls
 while consumption-equivalent real income rises. Above
 \(\varepsilon+\alpha s\), both fall; equality with a reference line gives
 zero marginal change for the corresponding income measure. Open markers
 identify drag maxima, not optimal platform shares. The regions classify local
 marginal effects, not finite changes. Illustration, not calibration.}
 \label{fig:local-marginal-regions}
\end{figure}

\subsection{Boundary cases and outside-funded spending}

If \(\Delta(\kappa)=0\), then \(\Lambda_g\equiv0\), so no unique drag peak
exists. If \(\Delta(\kappa)<0\), integration shifts expenditure toward the
higher-retention channel and the capture-drag proposition does not apply. The
capacity derivatives in Proposition~\ref{prop:general-share-speed} are ordinary
interior derivatives; feasible one-sided versions apply at \(\kappa=0,1\).
Without symmetry of \(g\), the unique peak remains but need not lie below
one-half.

The endpoint-crossing condition in Proposition~\ref{prop:finite-capacity} is
sufficient for two interior capacity thresholds. A target strictly outside
\([\mathcal D(1),\mathcal D(0)]\) has no threshold on \([0,1]\); equality with
an endpoint places the threshold at \(\kappa=1\) or \(0\), respectively. Only a
target in \((\mathcal D(1),\mathcal D(0))\) yields an interior threshold. The
finite ratio identities remain valid in every case.

For completeness, when nominal outside-funded spending \(G>0\) is fixed,
\begin{align}
 \frac{d\ln Y}{dz}
 &=\frac{\varepsilon B+\omega_zG}{B+\omega G}
 +\frac{\alpha\omega_z}{D}, \\
 \frac{d\ln R}{dz}
 &=\frac{\varepsilon B+\omega_zG}{B+\omega G}
 +\frac{\alpha\omega_z}{D}+\alpha s, \\
 \frac{\partial\ln Y}{\partial\kappa}
 &=s\delta_\rho\left[\frac{G}{B+\omega G}+\frac{\alpha}{D}\right]>0,
 \qquad
 \frac{\partial\ln R}{\partial\kappa}
 =\frac{\partial\ln Y}{\partial\kappa}.
 \label{eq:g-positive-derivatives}
\end{align}
For a finite change, define
\[
 \mathcal N=\frac{B_1+\omega_1G}{B_0+\omega_0G}.
\]
Then \(Y_1/Y_0=\mathcal N/\mathcal D\) and
\(R_1/R_0=\mathcal N/(\mathcal D\mathcal P^\alpha)\). Because
\(\mathcal N\) itself varies with capacity, the closed-form M1 capacity
thresholds are stated for \(G=0\).
